\subsection{Spectral identification of the effective Hamiltonian}
  \label{subsec:schrodinger-hamiltonian}

  The previous subsection established that the restriction of the linearized relaxation
  operator to the projected-stable sector is anti-Hermitian,
  $\mathcal{L}|_{\mathrm{proj}} = -i\,\Omega$, with $\Omega$ Hermitian and real spectrum
  $\{\omega_n\}$.
  The present subsection closes the remaining open step: identifying $\Omega$ with a
  function of the substrate Laplacian, and deriving the standard form of $\hat{H}_{\mathrm{eff}}$
  as a consequence.

  \paragraph{Step A: Linearized dispersion in the spectrally admissible regime.}

    In the spectrally admissible regime, the substrate modes $u_n$ satisfy
    $\mathcal{L}\,u_n = -i\omega_n u_n$ with $\omega_n \in \mathbb{R}$.
    The same modes satisfy the relational Laplacian eigenvalue equation
    $\Delta_\chi\,u_n = \lambda_n\,u_n$,
    where $\Delta_\chi$ denotes the substrate Laplacian associated with the relational
    connectivity of $\chi$.
    In the linearized (weak-amplitude) regime, the dispersion relation between the
    effective phase frequency $\omega_n$ and the Laplacian eigenvalue $\lambda_n$ follows
    from the admissibility envelope derived in~\cite{Beau2026a1}:
    \[
      \omega_n^2 = \lambda_n + O\!\left(\lambda_n^2 A_n^2 / c_\chi^2\right),
    \]
    where the correction term is controlled by the Born--Infeld saturation parameter
    $A_n^2/c_\chi^2 \ll 1$ in the weak-amplitude regime.
    To leading order one therefore has
    \[
      \omega_n = \sqrt{\lambda_n},
    \]
    so that $\Omega = \sqrt{\Delta_\chi}|_{\mathrm{proj}}$ in the projected-stable sector.

  \paragraph{Step B: Projected Laplacian and effective metric.}

    The projection $\Pi$ restricts the substrate Laplacian to the admissible spectral
    subspace via
    \[
      \Delta_\Pi = P^\dagger \Delta_\chi\, P,
    \]
    where $P$ denotes the spectral projection operator onto the projected-stable sector
    (Section~5.1).
    In the continuum limit of a locally homogeneous and isotropic relational substrate,
    the principal symbol of $\Delta_\chi$ defines an effective metric tensor $g^{\mu\nu}(x)$
    through the operator--metric correspondence established in~\cite{Beau2026c}:
    \[
      \sigma_2(\Delta_\chi)(x,k) = g^{\mu\nu}(x)\,k_\mu k_\nu.
    \]
    The projected Laplacian $\Delta_\Pi$ therefore inherits this metric structure,
    and in the low-lying spectral sector one has
    \[
      \Delta_\Pi \simeq \Delta_g + O(\ell_\chi^2\,\Delta_g^2),
    \]
    where $\Delta_g$ is the scalar Laplace--Beltrami operator of the effective metric $g$,
    and corrections are suppressed by the pre-geometric scale $\ell_\chi$.

  \paragraph{Step C: Effective Hamiltonian from spectral identification.}

    Combining Steps A and B, the effective phase frequencies satisfy
    \[
      E_n = \hbar_{\mathrm{eff}}\,\omega_n = \hbar_{\mathrm{eff}}\,\sqrt{\lambda_n},
    \]
    with $\lambda_n$ the eigenvalues of $\Delta_\Pi \simeq \Delta_g$.
    The effective Hamiltonian $\hat{H}_{\mathrm{eff}}$ defined spectrally by
    $\hat{H}_{\mathrm{eff}}\phi_n = E_n\phi_n$ therefore acts on projected spatial modes as
    \[
      \hat{H}_{\mathrm{eff}} = \hbar_{\mathrm{eff}}\,\sqrt{\Delta_g}
      + O(\ell_\chi^2\,\Delta_g^{3/2}).
    \]
    In the non-relativistic regime, one restricts to low-lying modes for which
    $\lambda_n \ll \ell_\chi^{-2}$.
    Expanding the square root to second order in $k$ around $k = 0$ gives
    \[
      \hbar_{\mathrm{eff}}\,\sqrt{\Delta_g}
      \simeq \hbar_{\mathrm{eff}}\left(m_{\mathrm{eff}} c^2
      + \frac{\Delta_g}{2m_{\mathrm{eff}}} + \cdots\right),
    \]
    where $m_{\mathrm{eff}}$ is the inertial parameter of the mode family, defined by
    the spectral gap separating the zero mode from the first admissible excited mode.
    Subtracting the rest-energy contribution and working in natural units for the
    effective description, one recovers
    \[
      \hat{H}_{\mathrm{eff}}
      = -\frac{\hbar_{\mathrm{eff}}^2}{2m_{\mathrm{eff}}}\,\Delta_g + V_{\mathrm{eff}}(x)
      + O(\ell_\chi^2\,\Delta_g^2),
    \]
    where $V_{\mathrm{eff}}(x)$ encodes the projected potential generated by background
    relational inhomogeneities.
    This is the standard non-relativistic Schr\"{o}dinger Hamiltonian, derived here as
    the leading-order effective operator of the projected substrate dynamics.

  \paragraph{Structural summary and scope.}

    The complete derivation chain is:
    \[
      \partial_\tau \chi = \mathcal{R}[\chi]
      \;\xrightarrow{\text{linearize}}\;
      \mathcal{L}|_{\mathrm{proj}} = -i\,\Omega
      \;\xrightarrow{\omega_n^2 = \lambda_n}\;
      \Omega = \sqrt{\Delta_\Pi}
      \;\xrightarrow{\Delta_\Pi \simeq \Delta_g}\;
      \hat{H}_{\mathrm{eff}} \simeq -\frac{\hbar_{\mathrm{eff}}^2}{2m_{\mathrm{eff}}}\,\Delta_g.
    \]

    Three limitations must be stated explicitly.

    First, the identification $\omega_n^2 = \lambda_n$ holds in the linearized regime.
    The Born--Infeld correction modifies this as
    $\omega_n^2(A_n) \approx \lambda_n(1 - 3\lambda_n A_n^2/8c_\chi^2)$,
    inducing an amplitude-dependent softening of the effective Hamiltonian that is
    negligible for $A_n \ll A_n^{\max}$ but becomes significant near saturation.

    Second, the operator--metric correspondence $\sigma_2(\Delta_\chi) = g^{\mu\nu}k_\mu k_\nu$
    is established in the continuum limit of a locally homogeneous substrate.
    In strongly inhomogeneous or pre-geometric regimes, the principal symbol may depart
    from a metric tensor, and $\hat{H}_{\mathrm{eff}}$ receives non-metric corrections.

    Third, the inertial parameter $m_{\mathrm{eff}}$ is identified here through the spectral
    gap of the low-lying sector.
    Its precise numerical value for specific particle species requires a microscopic
    identification of the relevant mode family, which is addressed in companion work on
    the spectral mass spectrum~\cite{Beau2026e}.

    Within these stated limits, the derivation establishes that the standard
    non-relativistic quantum Hamiltonian is not an independent postulate.
    It is the leading-order effective operator of projected substrate dynamics,
    controlled by the substrate Laplacian, the effective metric, and the emergent
    projective quantum $\hbar_{\mathrm{eff}}$.
